\subsection{Mixed concave--convex residuals}\label{sec:cr-mixed}

Consider a rational residual problem
$\min\{q(z)=\frac12z^{\T}Az+b^{\T}z+c : z\in\prod_{i\in\mathcal R}[l_i,u_i]\}$,
for example the residual after the core has been fixed. Assume an
orientation $o$ with $o_io_jA_{ij}\le0$ for $i\ne j$, and a partition
$\mathcal R=\mathcal D\cup\mathcal P$ with $A_{ii}\le0$ for $i\in\mathcal D$ and
$A_{\mathcal P\mathcal P}\succeq0$. Coordinates in $\mathcal P$ are
continuous; coordinates in $\mathcal D$ may be integer with inward-rounded
bounds. Define $\alpha_i,d_i$ as before and, for
$S\subseteq\mathcal D$, the label $\zeta(S)_i=\alpha_i+d_i[i\in S]$ and
\[
g(S)=\min\Bigl\{q(\zeta(S),z_{\mathcal P}) :
z_{\mathcal P}\in\prod_{i\in\mathcal P}[l_i,u_i]\Bigr\}.
\]
With $\mathcal P=\emptyset$ this is the cut class; with
$\mathcal D=\emptyset$ it is convex.

\begin{proposition}[Submodular conditional value]\label{prop:cr-submod}
(a) The minimum of $q$ over the continuous box and over the mixed domain
equals $\min_{S\subseteq\mathcal D}g(S)$.
(b) Each $g(S)$ is the value of a rational convex box QP.
(c) For all $S,T\subseteq\mathcal D$,
$g(S)+g(T)\ge g(S\cap T)+g(S\cup T)$.
\end{proposition}
\begin{proof}
(a) As in Lemma~\ref{lem:cr-endpoint}, move each coordinate of
$\mathcal D$ to an endpoint without increasing $q$. (b) After fixing
$z_{\mathcal D}$, the objective in $z_{\mathcal P}$ has Hessian
$A_{\mathcal P\mathcal P}\succeq0$. (c) Substitute out degenerate intervals
and put $z_i=\alpha_i+d_it_i$, $t\in[0,1]^{\mathcal R}$; the transformed objective
$\tilde q$ has off-diagonal Hessian entries
$d_id_jA_{ij}=o_io_j(u_i-l_i)(u_j-l_j)A_{ij}\le0$. For $s,t\in[0,1]^{\mathcal R}$,
$\tilde q(s)+\tilde q(t)\ge\tilde q(s\wedge t)+\tilde q(s\vee t)$: unary terms
cancel, and for a pair $i,j$ the two sides agree if $s_i-t_i$ and
$s_j-t_j$ have the same sign, while otherwise
$s_is_j+t_it_j-(s\wedge t)_i(s\wedge t)_j-(s\vee t)_i(s\vee t)_j
=(s_i-t_i)(s_j-t_j)\le0$, which multiplied by a nonpositive coefficient is
nonnegative. Let $t^S,t^T$ attain $g(S),g(T)$ in transformed
coordinates. Then $(\mathbf 1_S,t^S)\wedge(\mathbf 1_T,t^T)
=(\mathbf 1_{S\cap T},t^S\wedge t^T)$ and similarly for $\vee$ with
$S\cup T$; both completions are feasible. Hence
$g(S)+g(T)\ge\tilde q(\mathbf 1_{S\cap T},t^S\wedge t^T)
+\tilde q(\mathbf 1_{S\cup T},t^S\vee t^T)\ge g(S\cap T)+g(S\cup T)$.
\end{proof}

Positive semidefiniteness is not used in (c); it makes each $g(S)$
computable exactly in polynomial time \cite{KozlovTarasovKhachiyan1980}.
Partial minimization preserving submodularity is classical
\cite{Topkis1978}, and the same composition with submodular minimization
appears in \cite{BuntonTabuada2022,GomezHan2025}.

Let $m=|\mathcal D|$ and $h(S)=g(S)-g(\emptyset)$. For an ordering $\pi$ of
$\mathcal D$ with prefixes $S_0=\emptyset,S_1,\ldots,S_m=\mathcal D$, the
greedy vector is $b^\pi_{\pi(t)}=h(S_t)-h(S_{t-1})$.

\begin{proposition}[Greedy-base certificate]\label{prop:cr-greedy}
(a) For every ordering $\pi$ and every $T\subseteq\mathcal D$,
$b^\pi(T)\le h(T)$ and $b^\pi(\mathcal D)=h(\mathcal D)$.
(b) If $w=\sum_\pi\lambda_\pi b^\pi$ is a convex combination, then
$\min q\ge g(\emptyset)+\sum_{i\in\mathcal D}\min\{0,w_i\}$.
(c) Some convex combination of at most $m+1$ greedy vectors attains
equality \cite{Edmonds1970}, and such a combination, a minimizing label
and exact convex-QP witnesses can be computed in time polynomial in the
input length \cite{GrotschelLovaszSchrijver1988}.
(d) Given orderings, weights, the $O(m^2)$ chain values with their
minimizers, and a label $S^*$ with its completion, a checker verifies the
bound in (b) and its attainment by $S^*$ with polynomially many rational
operations and no optimization. It checks a convex chain value by
feasibility, exact evaluation and the sign conditions
$\partial_iq=0$ at interior, $\ge0$ at lower and $\le0$ at upper bounds
of $\mathcal P$, which suffice by convexity once
$A_{\mathcal P\mathcal P}\succeq0$ has been verified exactly.
\end{proposition}
\begin{proof}
(a) List $T=\{t_1,\ldots,t_r\}$ in $\pi$-order and let $S^{(s)}$ be the
$\pi$-prefix ending just before $t_s$. Since
$\{t_1,\ldots,t_{s-1}\}\subseteq S^{(s)}$, submodularity gives
$b^\pi_{t_s}=h(S^{(s)}\cup\{t_s\})-h(S^{(s)})
\le h(\{t_1,\ldots,t_s\})-h(\{t_1,\ldots,t_{s-1}\})$. Summing telescopes
to $b^\pi(T)\le h(T)$; the full prefix sum is $h(\mathcal D)$.
(b) By (a) and convexity, $w(S)\le h(S)$ for every $S$, so
$g(S)=g(\emptyset)+h(S)\ge g(\emptyset)+w(S)\ge g(\emptyset)+\sum_i\min\{0,w_i\}$;
conclude with Proposition~\ref{prop:cr-submod}(a).
(c) Edmonds' theorem gives
$\min_Sh(S)=\max\{\sum_i\min\{0,w_i\}: w\in B(h)\}$, where $B(h)$ is the
base polytope; its vertices are the greedy vectors, and the maximum of
this concave piecewise-linear function is attained at a combination of at
most $m+1$ of them. For computation, consider the linear program
$\max\{\beta: 0\le y\le\mathbf 1,\ \beta+(b^\pi)^{\T}y\le0\ \forall\pi\}$,
whose dual is $\min\{\sum_ir_i : \sum_\pi\lambda_\pi b^\pi+r\ge0,\
\sum_\pi\lambda_\pi=1,\ \lambda,r\ge0\}$. Sorting $y$ decreasingly and
forming its greedy vector maximizes $(b^\pi)^{\T}y$, so separation
costs $m+1$ convex-QP queries. Each $g(S)$ is the optimal value of a convex
box QP with fixed Hessian, and by the smallest-face argument of
Lemma~\ref{lem:cr-height} its encoding length is polynomial; hence the
polyhedron has polynomial facet complexity, and the ellipsoid method
solves the program and recovers an optimal basic dual solution from the
separating inequalities in polynomial time
\cite[Thm.~6.4.9, Lemma~6.5.15]{GrotschelLovaszSchrijver1988}. A basic dual
solution has at most $m+1$ positive $\lambda_\pi$, and at optimality
$r_i=\max\{0,-w_i\}$. For an optimal $y$ sorted as
$y_{\pi(1)}\ge\cdots\ge y_{\pi(m)}$, with $y_{\pi(m+1)}:=0$,
$(b^\pi)^{\T}y=\sum_t(y_{\pi(t)}-y_{\pi(t+1)})h(S_t)+(1-y_{\pi(1)})h(\emptyset)$
is a convex combination of prefix values equal to $\min h$, so some prefix
with positive weight is a minimizing label. (d) follows from (a), (b) and
the stated optimality conditions for convex problems on a box.
\end{proof}
